\section{Silhouette coefficient}
\label{sec:ex7}
\begin{problem}
For two observations, the average within-cluster distance $a(i)$ and nearest other-cluster average distance $b(i)$ are:
\begin{center}\begin{tabular}{@{}lrr@{}}\toprule Observation&$a(i)$&$b(i)$\\\midrule $x_1$&3&9\\$x_2$&5&4\\\bottomrule\end{tabular}\end{center}
Use $s(i)=[b(i)-a(i)]/\max\{a(i),b(i)\}$.
\begin{enumerate}[label=\textbf{(\alph*)}]
\item Calculate both silhouette coefficients.
\item Interpret the sign of each coefficient.
\item Does a high average silhouette guarantee meaningful domain-specific groups? Explain.
\end{enumerate}
\end{problem}


\subq{(a) Calculate both silhouettes}
Use the exact formula supplied by the guide:
\[
s(i)=\frac{b(i)-a(i)}{\max\{a(i),b(i)\}}.
\]
For $x_1$, $b=9>a=3$, so the denominator is $9$:
\[
s(x_1)=\frac{9-3}{\max\{3,9\}}=\frac{6}{9}
=\boxed{\frac{2}{3}\approx0.667}.
\]
For $x_2$, $b=4<a=5$, so the denominator is $5$:
\[
s(x_2)=\frac{4-5}{\max\{5,4\}}=-\frac15
=\boxed{-0.2}.
\]

\subq{(b) Interpret the signs}
$x_1$ has a positive silhouette, suggesting that on average it is substantially closer to its own cluster than to its closest competing cluster. $x_2$ has a negative silhouette: on average it lies closer to an alternative cluster than to its assigned one, potentially indicating questionable placement or overlapping groups. A coefficient near $0$ indicates a borderline observation; the range is $[-1,1]$.

\subq{(c) Does a high mean silhouette guarantee domain meaning?}
No. Silhouette quantifies within- versus between-cluster \emph{distance} relationships; it does not test semantic validity, causal explanation, external labels, clinical usefulness, or interpretability. For example, document clusters might separate by document length or writing style while failing to correspond to desired research topics. Domain review, cluster exemplars, stability, and external validation (if available) remain important.
\begin{answer}
$\boxed{s(x_1)\approx0.667}$ indicates good separation; $\boxed{s(x_2)=-0.2}$ suggests possible misassignment or overlap.
\end{answer}


\subsection*{Supplementary worked example: calculate $a(i)$ and $b(i)$ from raw observations}
\begin{reviewproblem}
Consider three \emph{given} clusters on the real number line, with ordinary Euclidean distance
$d(x,y)=|x-y|$:
\[
C_1=\{0,2,7\},\qquad C_2=\{8,9\},\qquad C_3=\{15,18\}.
\]
\begin{enumerate}[label=\textbf{(\alph*)}]
\item Calculate the within-cluster mean distance $a(i)$ and the nearest other-cluster mean distance $b(i)$ for $x_i=2$.
\item Repeat for $x_i=7$ and calculate the silhouette of each observation.
\item Explain why $b(i)$ is \emph{not} simply the distance to the nearest observation in another cluster.
\end{enumerate}
These numbers are a constructed practice example; they are \emph{not} additional observations supplied by the original study guide.
\end{reviewproblem}

\begin{center}
\begin{tikzpicture}[x=0.44cm,y=1cm]
  \draw[-{Latex[length=2mm]},gray!75,thick] (-0.7,0)--(19.7,0) node[right] {$x$};
  \foreach \x in {0,2,7,8,9,15,18}{\draw[gray] (\x,.09)--(\x,-.09);}
  \foreach \x in {0,2,7}{\fill[navy] (\x,0) circle (2.5pt);}
  \foreach \x in {8,9}{\fill[forest] (\x,0) circle (2.5pt);}
  \foreach \x in {15,18}{\fill[amber] (\x,0) circle (2.5pt);}
  \foreach \x in {0,2,7,8,9,15,18}{\node[below=3pt,font=\scriptsize] at (\x,0) {$\x$};}
  \node[above=8pt,text=navy,font=\small] at (2,0) {$C_1$};
  \node[above=8pt,text=forest,font=\small] at (8.5,0) {$C_2$};
  \node[above=8pt,text=amber,font=\small] at (16.5,0) {$C_3$};
\end{tikzpicture}
\end{center}

\textbf{Definitions.} If $C(i)$ is the assigned cluster containing $x_i$, then
\[
\begin{aligned}
a(i)&=\frac{\sum_{x_j\in C(i),\,j\neq i}d(x_i,x_j)}{|C(i)|-1},\\
b(i)&=\min_{C\neq C(i)}\underbrace{\left[\frac{1}{|C|}\sum_{x_j\in C}d(x_i,x_j)\right]}_{\text{mean distance to an entire alternative cluster}}.
\end{aligned}
\]
The observation itself is \emph{excluded} when computing $a(i)$; \emph{every} member of each competing cluster is included when computing $b(i)$.

\subq{(a) Point $x_i=2$, assigned to $C_1$}
\textbf{Step 1 --- its own cluster.} The other members of $C_1$ are $0$ and $7$, hence
\[
a(2)=\frac{|2-0|+|2-7|}{3-1}=\frac{2+5}{2}=\boxed{3.5}.
\]
\textbf{Step 2 --- each alternative cluster.} Do \emph{not} compare only the nearest points:
\[
\begin{aligned}
\bar d(2,C_2)&=\frac{|2-8|+|2-9|}{2}=\frac{6+7}{2}=6.5,\\
\bar d(2,C_3)&=\frac{|2-15|+|2-18|}{2}=\frac{13+16}{2}=14.5.
\end{aligned}
\]
\textbf{Step 3 --- nearest alternative cluster in average distance.}
\[
b(2)=\min\{6.5,14.5\}=\boxed{6.5},\qquad
s(2)=\frac{6.5-3.5}{\max\{3.5,6.5\}}=\boxed{\frac{6}{13}\approx0.462}.
\]

\subq{(b) Point $x_i=7$, also assigned to $C_1$}
\textbf{Step 1.} Excluding $7$ itself gives
\[
a(7)=\frac{|7-0|+|7-2|}{2}=\frac{7+5}{2}=\boxed{6}.
\]
\textbf{Step 2.} Average distance to each other cluster:
\[
\bar d(7,C_2)=\frac{|7-8|+|7-9|}{2}=\frac{1+2}{2}=1.5,
\]
\[
\bar d(7,C_3)=\frac{|7-15|+|7-18|}{2}=\frac{8+11}{2}=9.5.
\]
\textbf{Step 3.} Therefore
\[
b(7)=\min\{1.5,9.5\}=\boxed{1.5},\qquad
s(7)=\frac{1.5-6}{\max\{6,1.5\}}=\boxed{-0.75}.
\]
A negative silhouette means that $7$ is on average \emph{closer to cluster $C_2$} than to the other observations of its assigned cluster $C_1$. This is a sign of possible misassignment, not proof that the clustering algorithm made a computational error.

\subq{(c) Why the nearest individual point is not $b(i)$}
For $x_i=2$, the nearest observation outside $C_1$ is $8$, at distance $6$. But $b(2)=6.5$ because it is the \emph{mean of the two distances} $6$ and $7$ to the whole closest competing cluster. With more than two clusters, evaluate each competing cluster's \emph{average} first, then select the smallest average.

\begin{answer}
\textbf{Memory rule:} $a(i)$ = average distance to \emph{other observations in your own cluster}; $b(i)$ = smallest \emph{cluster-average} distance among all competing clusters. Here,
\[
\boxed{(a(2),b(2),s(2))=(3.5,6.5,6/13)},\qquad
\boxed{(a(7),b(7),s(7))=(6,1.5,-0.75)}.
\]
\end{answer}
